%% -*- coding: utf-8 -*-

\section{被动语态及虚拟语气}

\subsection{一般时（现在、过去、将来）的被动语态}

\begin{center}
  \begin{talltblr}[
    label = none, % Removes the "Table X:" text
    ]{
    vlines, hlines,
    % 加入支持列表环境的测量参数 (记得在导言区加载 \UseTblrLibrary{varwidth})
    % measure=vbox, stretch=-1,
    rows={halign=c, valign=m},
    columns={halign=c, valign=m},
    % column{1}={0.15\linewidth, halign=c, bg=azure9},
    % column{2}={0.35\linewidth, halign=l},
    % column{3}={0.35\linewidth, halign=l},
    column{4}={0.25\linewidth, halign=l},
    row{1}={font=\large\bfseries, bg=azure9, halign=c}
    }
    时态&句式&构成&举例\\
    \SetCell[r=4]{c}一般现在时&肯定句&主语+am/is/are done&\SetCell[r=4]{l}The decision is made by the board.\\
        &否定句&主语+am/is/are not done&\\
        &一般疑问句&am/is/are+主语+done?&\\
        &特殊疑问句&特殊疑问词+am/is/are+主语+done?&\\
    \SetCell[r=4]{c}一般过去时&肯定句&主语+was/were done&\SetCell[r=4]{l}We were not invited to the event.\\
        &否定句&主语+was/were not done&\\
        &一般疑问句&was/were+主语+done?&\\
        &特殊疑问句&特殊疑问词+was/were+主语+done?&\\
    \SetCell[r=4]{c}一般将来时&肯定句&主语+will be done&\SetCell[r=4]{l}When will the contract be renewed?\\
        &否定句&主语+will not be done&\\
        &一般疑问句&Will+主语+be done&\\
        &特殊疑问句&特殊疑问词+will+主语+be done?&\\
  \end{talltblr}
\end{center}

\subsection{进行时（现在、过去）的被动语态}

\begin{center}
  \begin{talltblr}[
    label = none, % Removes the "Table X:" text
    ]{
    vlines, hlines,
    % 加入支持列表环境的测量参数 (记得在导言区加载 \UseTblrLibrary{varwidth})
    % measure=vbox, stretch=-1,
    rows={halign=c, valign=m},
    columns={halign=c, valign=m},
    % column{1}={0.15\linewidth, halign=c, bg=azure9},
    % column{2}={0.35\linewidth, halign=l},
    % column{3}={0.35\linewidth, halign=l},
    column{4}={0.25\linewidth, halign=l},
    row{1}={font=\large\bfseries, bg=azure9, halign=c}
    }
    时态&句式&构成&举例\\
    \SetCell[r=4]{c}现在进行时&肯定句&主语+am/is/are being done&\SetCell[r=4]{l}The data is being analyzed by the researchers.\\
        &否定句&主语+am/is/are not being done&\\
        &一般疑问句&am/is/are+主语+being done?&\\
        &特殊疑问句&特殊疑问词+am/is/are+主语+being done?&\\
    \SetCell[r=4]{c}过去进行时&肯定句&主语+was/were being done&\SetCell[r=4]{l}The room was being cleaned by him at nine yesterday morning.\\
        &否定句&主语+was/were not being done&\\
        &一般疑问句&was/were+主语+being done?&\\
        &特殊疑问句&特殊疑问词+was/were+主语+being done?&\\
  \end{talltblr}
\end{center}

\subsection{完成时（现在、过去）的被动语态}

\begin{center}
  \begin{talltblr}[
    label = none, % Removes the "Table X:" text
    ]{
    vlines, hlines,
    % 加入支持列表环境的测量参数 (记得在导言区加载 \UseTblrLibrary{varwidth})
    % measure=vbox, stretch=-1,
    rows={halign=c, valign=m},
    columns={halign=c, valign=m},
    % column{1}={0.15\linewidth, halign=c, bg=azure9},
    % column{2}={0.35\linewidth, halign=l},
    % column{3}={0.35\linewidth, halign=l},
    column{4}={0.25\linewidth, halign=l},
    row{1}={font=\large\bfseries, bg=azure9, halign=c}
    }
    时态&句式&构成&举例\\
    \SetCell[r=4]{c}现在完成时&肯定句&主语+have/has been done&\SetCell[r=4]{l}The project has not been completed yet.\\
        &否定句&主语+have/has not been done&\\
        &一般疑问句&Have/Has+主语+been done?&\\
        &特殊疑问句&特殊疑问词+have/has+主语+been done?&\\
    \SetCell[r=4]{c}过去完成时&肯定句&主语+had been done&\SetCell[r=4]{l}Had the road been blocked when you passed by?\\
        &否定句&主语+had been done&\\
        &一般疑问句&Had+主语+been done?&\\
        &特殊疑问句&特殊疑问词+had+主语+being been done?&\\
  \end{talltblr}
\end{center}

\subsection{情态动词的被动语态}

\begin{center}
  \begin{talltblr}[
    label = none, % Removes the "Table X:" text
    ]{
    vlines, hlines,
    % 加入支持列表环境的测量参数 (记得在导言区加载 \UseTblrLibrary{varwidth})
    % measure=vbox, stretch=-1,
    rows={halign=c, valign=m},
    columns={halign=c, valign=m},
    % column{1}={0.15\linewidth, halign=c, bg=azure9},
    % column{2}={0.35\linewidth, halign=l},
    % column{3}={0.35\linewidth, halign=l},
    column{4}={0.25\linewidth, halign=l},
    row{1}={font=\large\bfseries, bg=azure9, halign=c}
    }
    句式&构成&举例\\
    肯定句&主语+情态动词+be done+其他&The rules must be followed by every student.\\
    否定句&主语+情态动词+not be done+其他&The decision needn't be made immediately.\\
    一般疑问句&情态动词+主语+be done+其他?&Can the work be finished ahead of schedule?\\
    特殊疑问句&特殊疑问词+情态动词+主语+be done+其他?&Where could the lost keys be found?\\
  \end{talltblr}
\end{center}

\subsection{常用被动语态的短语}

\begin{enumerate}
  \item \textbf{常用被动语态的短语}： be appointed as被任命为…、be blamed for因为……被责备、be composed of由……组成、be equipped with装备有…、be known as 被称为…、be made of由…制成、be associated with 涉及;与……有关、be compared with与……相比、be decorated with用…装饰、be filled with 被充满…、be located in 坐落于…、be named after 以……命名
  \item \textbf{没有被动语态的短语}：break out 爆发、belong to属于……、go bad (食物等)变质、fall asleep 入睡、make progress 取得进步、keep in touch with与……保持联系、set off出发，动身、pass away 去世、take place发生、turn up 出现
\end{enumerate}

\subsection{动词主动形式表示被动含义}

\begin{center}
  \begin{talltblr}[
    label = none, % Removes the "Table X:" text
    ]{
    vlines, hlines,
    % 加入支持列表环境的测量参数 (记得在导言区加载 \UseTblrLibrary{varwidth})
    % measure=vbox, stretch=-1,
    rows={halign=c, valign=m},
    columns={halign=c, valign=m},
    column{1}={0.15\linewidth, halign=c, bg=azure9},
    column{2}={0.35\linewidth, halign=l},
    column{3}={0.35\linewidth, halign=l},
    % column{4}={0.35\linewidth, halign=l},
    row{1}={font=\large\bfseries, bg=azure9, halign=c}
    }
    情况&例词&举例\\
    部分\(vi\).&change, open, develop, improve, stop, close等&His health improves gradually.\\
    表示感官的系动词&sound, taste, feel, look等&The plan sounds feasible.\\
    由well、easily、badly等修饰的部分\(vi\).&sell, read, write, wear&The latest smartphone sells extremely well.\\
  \end{talltblr}
\end{center}

\subsection{被动语态中to的还原}

\begin{enumerate}
  \item 在不定式作宾语补足语的结构中，句子变为被动语态时，原主动语态中的宾语提前至主语位置，宾语补足语(即不定式短语)的位置不变，比如:
        \begin{itemize}
          \item The teacher encouraged the students to participate in the competition.老师鼓励学生参加比赛。 (主动)
          \item The students were encouraged to participate in the competition.学生们被鼓励参加比赛。 (被动)
        \end{itemize}
  \item 不带to的不定式作宾语补足语，在变为被动语态时，省略的to要加上，常考的动词有make、let、have、see、watch、notice、hear 等，本句中的 made to 就是此用法。再比如:
        \begin{itemize}
          \item They heard her sing in the next room.他们听见她在隔壁房间唱歌。 (主动)
          \item She was heard to sing in the next room.有人听见她在隔壁房间唱歌。 (被动)
        \end{itemize}
\end{enumerate}

\subsection{其他表示被动语态的方式}

\begin{enumerate}
  \item 被动语态除了用“be+过去分词”来表示外，还可以用以下结构来表示。
        \begin{center}
          \begin{talltblr}[
            label = none, % Removes the "Table X:" text
            ]{
            vlines, hlines,
            % 加入支持列表环境的测量参数 (记得在导言区加载 \UseTblrLibrary{varwidth})
            % measure=vbox, stretch=-1,
            rows={halign=c, valign=m},
            columns={halign=c, valign=m},
            % column{1}={0.15\linewidth, halign=c, bg=azure9},
            column{1}={0.25\linewidth, halign=l},
            column{2}={0.55\linewidth, halign=l},
            % column{4}={0.35\linewidth, halign=l},
            row{1}={font=\large\bfseries, bg=azure9, halign=c}
            }
            结构&举例\\
            get+过去分词&The window got broken during the storm.窗户在暴雨中被打破了。\\
            be+形容词+不定式&The problem is hard to solve.这个问题很难被解决。\\
            make+形容词+不定式&The teacher made the lesson easy to understand.老师把这节课讲得通俗易懂。\\
            need+动名词&The house needs painting.这栋房子需要被粉刷。\\
          \end{talltblr}
        \end{center}
  \item 使用“be+过去分词”和“get+过去分词”表示被动语态时，二者之间有一定区别。
        \begin{itemize}
          \item “be+过去分词”多用来表示一般的动作或情况;
          \item “get+过去分词”多用于表示动作的结果或动作的渐变。
        \end{itemize}
  \item “get+过去分词”结构变疑问句或否定句时，需要添加助动词do或did，比如:How did that blue door get locked?那扇蓝色的门是怎么锁上的?
\end{enumerate}

\subsection{被动语态的固定用法}

\begin{enumerate}
  \item It+被动态动词+that
        \begin{itemize}
          \item It is said that\(\ldots\) 据说…
          \item It is well known that\(\ldots\) 众所周知…
          \item It is believed that\(\ldots\) 人们相信…
          \item It is hoped that\(\ldots\) 人们希望…
          \item It is found that \(\ldots\) 人们发现…
          \item It is reported that\(\ldots\) 据报道…
        \end{itemize}
  \item 主语+被动态动词+不定式： 主语+ be asked / be believed / be considered / be estimated / be expected / be felt / be known / be meant / be reported / be said / be seen / be supposed / be thought /be understood+不定式
\end{enumerate}

\subsection{被动语态的使用场景}

\begin{center}
  \begin{talltblr}[
    label = none, % Removes the "Table X:" text
    ]{
    vlines, hlines,
    % 加入支持列表环境的测量参数 (记得在导言区加载 \UseTblrLibrary{varwidth})
    % measure=vbox, stretch=-1,
    rows={halign=c, valign=m},
    columns={halign=c, valign=m},
    % column{1}={0.15\linewidth, halign=c, bg=azure9},
    column{2}={0.55\linewidth, halign=l},
    % column{3}={0.35\linewidth, halign=l},
    % column{4}={0.35\linewidth, halign=l},
    row{1}={font=\large\bfseries, bg=azure9, halign=c}
    }
    情况&举例\\
    不知道或不必提及动作的执行者&The house was robbed last night.房子昨晚被盗了。\\
    强调动作的承受者&The statue was carved by Michelangelo.这座雕像由米开朗基罗雕刻而成。\\
    动作的执行者是无生命的事物&The power was cut off by the lightning.电力被闪电切断了。\\
    用于陈述客观事实&Coffee is grown in many countries.咖啡在许多国家都有种植。\\
  \end{talltblr}
\end{center}

\subsection{虚拟条件句——与现在情况相反}

虚拟条件句可用于表达与实际情况相反的假设，以及在这种假设下可能产生的结果，常用于提出委婉的建议、表达遗憾或描述理想化情境，其结构如下:
\begin{center}
  \begin{talltblr}[
    label = none, % Removes the "Table X:" text
    ]{
    vlines, hlines,
    % 加入支持列表环境的测量参数 (记得在导言区加载 \UseTblrLibrary{varwidth})
    % measure=vbox, stretch=-1,
    rows={halign=c, valign=m},
    columns={halign=c, valign=m},
    % column{1}={0.15\linewidth, halign=c, bg=azure9},
    column{2}={0.29\linewidth, halign=l},
    column{3}={0.35\linewidth, halign=l},
    column{1}={0.25\linewidth, halign=l},
    row{1}={font=\large\bfseries, bg=azure9, halign=c}
    }
    从句&主句&举例\\
    If+主语+动词过去式 (be 动词要用 were)&主语+would/should/could/might+动词原形&If she were in your shoes, she could handle this problem easily.如果她处在你这种情况，她能轻松处理这个问题。\\
  \end{talltblr}
\end{center}

\subsection{虚拟条件句——与将来情况相反}

虚拟条件句可用来表达与将来情况相反的假设，以及在这种假设下可能产生的结果，常用于表达建议、描述愿望以及不太可能实现的情况，其结构如下:
\begin{center}
  \begin{talltblr}[
    label = none, % Removes the "Table X:" text
    ]{
    vlines, hlines,
    % 加入支持列表环境的测量参数 (记得在导言区加载 \UseTblrLibrary{varwidth})
    % measure=vbox, stretch=-1,
    rows={halign=c, valign=m},
    columns={halign=c, valign=m},
    % column{1}={0.15\linewidth, halign=c, bg=azure9},
    column{2}={0.29\linewidth, halign=l},
    column{3}={0.35\linewidth, halign=l},
    column{1}={0.25\linewidth, halign=l},
    row{1}={font=\large\bfseries, bg=azure9, halign=c}
    }
    从句&主句&举例\\
    {if+主语+动词过去式\\if+主语+should+动词原形\\if+主语+were to+动词原形}&主语+would/should/could/might+动词原形&If everyone were to reduce their carbon footprint, the environment would improve greatly.\\
  \end{talltblr}
\end{center}

\subsection{虚拟条件句——与过去情况相反}

虚拟条件句可用来表达与过去情况相反的假设，以及在这种假设下可能产生的结果，常用于表达遗憾、后悔以及对过去情况的推测，其结构如下:
\begin{center}
  \begin{talltblr}[
    label = none, % Removes the "Table X:" text
    ]{
    vlines, hlines,
    % 加入支持列表环境的测量参数 (记得在导言区加载 \UseTblrLibrary{varwidth})
    % measure=vbox, stretch=-1,
    rows={halign=c, valign=m},
    columns={halign=c, valign=m},
    % column{1}={0.15\linewidth, halign=c, bg=azure9},
    column{2}={0.29\linewidth, halign=l},
    column{3}={0.35\linewidth, halign=l},
    column{1}={0.25\linewidth, halign=l},
    row{1}={font=\large\bfseries, bg=azure9, halign=c}
    }
    从句&主句&举例\\
    if + 主语 + had + 动词过去分词&主语 + would/should/could/might + have + 过去分词&If they had double-checked the data, they would have avoided the error.\\
  \end{talltblr}
\end{center}

\subsection{would have done}

在虚拟语气中，would have done 用于非真实条件句中，表达与过去事实相反的假设，带有惋惜、遗憾、后悔等情感。比如:If I had studied harder, I would have passed the exam.如果我当时更努力学习，我就会通过考试了。类似的结构还有:
\begin{enumerate}
  \item \textbf{should have done}: 表示过去应该做某事但实际上没有做，含有责备或遗憾的意思。
        \begin{itemize}
          \item I should have gone to the doctor earlier. Now my cold is much worse.我本应该早点去看医生的。现在我的感冒严重了。
        \end{itemize}
  \item \textbf{could have done}：表示过去本来能够做某事但实际上没有做，强调能力或可能性缺失。
        \begin{itemize}
          \item I could have gone to the concert, but I didn't have the money.我本来可以去听音乐会，但我没有钱。
        \end{itemize}
\end{enumerate}

\subsection{wish用于表示虚拟语气}

“wish+宾语从句”表示不能或不太可能实现的愿望,后面的宾语从句经常使用虚拟语气,wish 可以理解为“要是……就好了”或者“但愿……”等，从句中谓语动词的时态要根据表达的时间来确定。
\begin{enumerate}
  \item 表示过去不能实现的愿望,从句的谓语动词用过去完成时 (had+动词过去分词)，比如
        \begin{itemize}
          \item I wish I had studied harder for the exam last week.我真希望上周我为考试更努力地学习了。
        \end{itemize}
  \item 表示现在不能实现的愿望，从句的谓语动词用过去式(be动词用were)，比如:
        \begin{itemize}
          \item We wish it were sunny today.我们真希望今天是晴天。
        \end{itemize}
  \item 表示将来不能实现的愿望，用“would(could)+动词原形”，比如:
        \begin{itemize}
          \item I wish our neighbours would stop making so much noise at night. 我真希望我们的邻居晚上不要再制造这么多噪声了。
        \end{itemize}
\end{enumerate}

\subsection{错综时间条件句}

当条件状语从句表示的时间和主句表示的时间不一致时，主句与从句谓语动词的形式要根据各自发生的时间来确定。
\begin{center}
  \begin{talltblr}[
    label = none, % Removes the "Table X:" text
    ]{
    vlines, hlines,
    % 加入支持列表环境的测量参数 (记得在导言区加载 \UseTblrLibrary{varwidth})
    % measure=vbox, stretch=-1,
    rows={halign=c, valign=m},
    columns={halign=c, valign=m},
    % column{1}={0.15\linewidth, halign=c, bg=azure9},
    column{2}={0.25\linewidth, halign=l},
    column{3}={0.45\linewidth, halign=l},
    % column{4}={0.35\linewidth, halign=l},
    row{1}={font=\large\bfseries, bg=azure9, halign=c}
    }
    从句&主句&举例\\
    \SetCell[r=2]{c}与过去试试修复&与现在或现在正在发生的事不符&If I had taken that job offer last year, I would have more financial stability now.如果我去年接受了那份工作邀请，我现在的经济状况会更稳定。\\
        &与未来结果不符&If she had saved enough money last year, she would be traveling around the world next month. 如果她去年存够了钱，下个月她就会在环游世界了。\\
    \SetCell[r=2]{c}与现在事实相反&与过去事实不符&If I were more outgoing, I would have made more friends at the party last night.如果我更外向一些，昨晚在派对上我就会交到更多朋友。\\
        &与未来结果不符&If I was at home, I'd be seeing my father tomorrow because it's his birthday.如果我在家，我明天就会去看我父亲，因为是他的生日。\\
    假设的未来情况&与现在或现在正在发生的事实不符&If I wasn't meeting my teacher later, I'd at home now.\\
  \end{talltblr}
\end{center}

\subsection{表示含蓄的虚拟语气}

在虚拟语气中并不总是会出现if条件句，这时就会通过其他手段来代替条件句，如本句中的 otherwise。其主要表现形式如下:
\begin{enumerate}
  \item 用 with、without、but for、otherwise、or、under 等代替条件状语从句，比如:
        \begin{itemize}
          \item But for his timely warning, we might have been in great danger.要不是他及时警告，我们可能已经陷人极大的危险之中了。
        \end{itemize}
  \item 条件暗含在上下文中，此时主、从句可以省略其中的一个，来表示说话人的一种强烈的感情，比如:
        \begin{itemize}
          \item You could have made a more reasonable decision. 你本可以做出一个更合理的决定。
        \end{itemize}
\end{enumerate}

\subsection{宾语从句中的虚拟语气}

\begin{enumerate}
  \item “wish+宾语从句”表示不能或不太可能实现的愿望时:
        \begin{enumerate}
          \item 用一般过去时表示与现在事实相反;
          \item 用过去完成时表示与过去事实相反;
          \item 用“would/should/could/might + 动词原形”表示与将来事实相反
        \end{enumerate}
        比如：
        \begin{itemize}
          \item He wishes he had apologized to his friend earlier. 他真希望当时自己早点向朋友道歉。
        \end{itemize}
  \item 在 would rather 后的宾语从句中:
        \begin{enumerate}
          \item 用一般过去时表示与现在或将来的事实相反;
          \item 用过去完成时表示与过去事实相反
        \end{enumerate}
        比如：
        \begin{itemize}
          \item He would rather she had told him the truth earlier. 他宁愿她当时早点告诉他真相。
        \end{itemize}
  \item 在表示建议、要求、命令等的动词后的宾语从句中也经常使用虚拟语气,宾语从句中谓语动词的结构为“should+动词原形”，其中should 经常省略，比如：
        \begin{itemize}
          \item The boss demanded that the project (should) be completed within a week. 老板要求这个项目在一周内完成。
        \end{itemize}
\end{enumerate}



%%% Local Variables:
%%% TeX-engine: xetex
%%% TeX-master: "../IELTS_300_Sentences"
%%% End:

% LocalWords:  vlines hlines halign valign xetex LocalWords
